\chapter{分布式计算工程实践}

分布式计算工程把算法、运行时、存储、网络和观测性结合起来。真正的系统不仅要在正常路径跑得快，还要在机器慢、网络抖动、任务失败和数据倾斜时保持可诊断、可恢复。

\section{Shuffle}

Shuffle 把数据按 key 重新分布，是很多分布式计算框架的核心成本。它包含序列化、压缩、网络传输、落盘、排序、合并和反序列化。一个 group-by 或 join 的成本，常常主要来自 shuffle，而不是用户函数。

优化 shuffle 的方向包括减少数据量、提前聚合、压缩、分区均衡、避免热点 key、使用本地磁盘顺序写和控制并发拉取。它是分布式版的数据布局问题。

\section{检查点和容错}

长任务必须考虑失败。检查点保存中间状态，让失败后从最近点恢复，而不是从头开始。检查点太频繁会拖慢正常路径，太稀疏会增加恢复成本。流式计算还要处理 exactly-once 或 at-least-once 语义。

\section{观测性}

没有观测性，分布式系统无法维护。日志、指标、trace、profile 和事件审计要能回答：请求经过哪些节点，在哪里等待，重试了几次，队列多长，哪个分片慢，哪个任务倾斜，哪次部署引入变化。

\section{从单机原则到集群原则}

第二册前面讲的原则在分布式系统中仍然成立。局部性对应数据本地调度；锁竞争对应热点 key；false sharing 对应多个任务争用同一分片；背压对应限流和队列控制；归约对应 map-side combine 和树形聚合；缓存一致性对应复制一致性和失效策略。系统规模变大，原则没有消失，只是成本单位从 cache line 变成网络消息和磁盘块。

\begin{keyidea}
分布式计算不是单机计算的替代，而是单机计算原则在通信、故障和复制约束下的扩展。
\end{keyidea}

\begin{labbox}
综合实验：设计一个小型分布式日志统计系统。输入按文件分片，worker 本地统计，shuffle 按 key 聚合，driver 合并结果。写清数据格式、失败重试、幂等输出、指标和瓶颈。
\end{labbox}
